%% ma: L12-B17-0002
%% lop: 12
%% chuong: 5
%% bai: 17
%% dang: 12.17.2
%% loai: TLN
%% muc: VD
%% dap_an: "3"
%% nguon: "0. ĐỀ THAM KHẢO TN THPT 2025.pdf (D00), Phần III Câu 3"
%% kien_thuc: "Bài 17 (mặt cầu, khoảng cách giữa hai điểm trong tọa độ); giải hệ bậc nhất ba ẩn bằng máy tính cầm tay"
%% trang_thai: cho-duyet
%% ly_do: "Dạng toán mới đề xuất 12.17.2, chờ giáo viên duyệt"
\begin{ex}
Hệ thống định vị toàn cầu GPS là một hệ thống cho phép xác định vị trí của một vật thể trong không gian. Trong cùng một thời điểm, vị trí của một điểm $M$ trong không gian sẽ được xác định bởi bốn vệ tinh cho trước nhờ các bộ thu phát tín hiệu đặt trên các vệ tinh. Giả sử trong không gian với hệ tọa độ $Oxyz$, có bốn vệ tinh lần lượt đặt tại các điểm $A(3;1;0)$, $B(3;6;6)$, $C(4;6;2)$, $D(6;2;14)$; vị trí $M(a;b;c)$ thỏa mãn $MA=3$, $MB=6$, $MC=5$, $MD=13$. Khoảng cách từ điểm $M$ đến điểm $O$ bằng bao nhiêu?
\shortans{3}
\loigiai{$M$ thuộc bốn mặt cầu tâm $A,B,C,D$. Khai triển $MA^2=9$, $MB^2=36$, $MC^2=25$, $MD^2=169$ rồi trừ vế theo vế phương trình đầu cho ba phương trình còn lại ta được hệ
\[\begin{cases}10b+12c=44\\ 2a+10b+4c=30\\ 6a+2b+28c=66\end{cases}\Leftrightarrow\begin{cases}a=1\\ b=2\\ c=2.\end{cases}\]
Kiểm tra $MA^2=4+1+4=9$. Vậy $M(1;2;2)$ và $OM=\sqrt{1+4+4}=3$.}
\end{ex}
